\section{บทนำ}
ข้อมูลการลานิสิตสามารถสะท้อนรูปแบบการขาดเรียนในแต่ละช่วงเวลา รายวิชา และกลุ่มนิสิต แต่ฐานข้อมูลที่รองรับงานประจำวันไม่ได้จัดรูปแบบมาเพื่อเปรียบเทียบหลายภาคการศึกษาโดยตรง คลังข้อมูลจึงทำหน้าที่รวบรวมบริบทที่จำเป็นและกำหนดหน่วยการวัดให้สอดคล้องกัน เพื่อให้ผู้ใช้แต่ละบทบาทมองข้อมูลจากรายละเอียดของตนไปถึงภาพรวมระดับคณะได้\cite{repo,warehouse}

วัตถุประสงค์ของบทความประกอบด้วยการอธิบายสถาปัตยกรรมคลังข้อมูล การกำหนดระดับรายละเอียดของข้อเท็จจริงและมิติ การเชื่อมข้อมูลเพื่อวิเคราะห์ และคำถามที่มุมมองแต่ละบทบาทตอบได้ ขอบเขตจำกัดอยู่ที่สถิติการลาและผลกระทบต่อการเรียน โดยไม่มีการนำภาพหน้าจอมาใช้ประกอบ

หลักฐานหลักคือโค้ดใน \nolinkurl{apps/api} และ \nolinkurl{apps/web} ของโครงการ ส่วนเอกสารสคีมาใช้ประกอบการอธิบายความสัมพันธ์ ไม่ถือว่าสคีมาที่ออกแบบหรือหน้ารายงานที่มีอยู่เป็นหลักฐานว่าข้อมูลจริงครบถ้วนหรือผ่านการประเมินประสิทธิภาพแล้ว
